\documentclass[11pt]{article}

% Packages
\usepackage[utf8]{inputenc}
\usepackage[T1]{fontenc}
\usepackage{amsmath,amssymb}
\usepackage{graphicx}
\usepackage[margin=1in]{geometry}
\usepackage{natbib}
\usepackage{hyperref}
\usepackage{setspace}
\usepackage{lineno}

\onehalfspacing
\linenumbers

\begin{document}

\section*{Introduction}

Neural activity at the mesoscopic scale---the coordinated fluctuations of local populations of neurons spanning multiple cortical regions---plays a central role in sensory processing, motor control, and cognition \citep{churchland2012neural, musall2019single, stringer2019high}. These dynamics unfold over spatial scales of millimeters and temporal scales of tens to hundreds of milliseconds, occupying a critical intermediate regime between the activity of single neurons and the global patterns observed in whole-brain imaging. Understanding how mesoscale cortical dynamics give rise to behavior remains a fundamental challenge in systems neuroscience.

Wide-field calcium imaging has emerged as a powerful tool for monitoring mesoscale cortical activity in behaving mice \citep{wekselblatt2016large, allen2017global, makino2017transformation}. By combining genetically encoded calcium indicators such as GCaMP with chronic cranial windows, wide-field imaging enables simultaneous recording of population-level fluorescence signals across the entire dorsal cortex at frame rates sufficient to resolve behaviorally relevant dynamics \citep{ren2021characterizing, ghanbari2019cortex}. Studies using this approach have revealed cortex-wide patterns of activation during sensorimotor tasks \citep{musall2019single, pinto2019task}, the emergence of coordinated cortical sequences during motor learning \citep{makino2017transformation}, and structured spontaneous activity that recapitulates functional connectivity motifs \citep{mohajerani2013spontaneous, vanni2014mesoscale}. These observations suggest that the spatiotemporal dynamics of mesoscale cortical activity carry rich information about the computational state of the brain.

A complementary line of work has sought to formalize neural population activity in terms of dynamical systems, in which the trajectory of population-level activity evolves according to rules that can be approximated mathematically. Seminal studies in motor cortex demonstrated that neural population dynamics during reaching could be captured by low-dimensional, approximately linear models \citep{churchland2012neural, cunningham2014dimensionality}. This dynamical systems perspective has since been applied across many brain areas and behavioral contexts, revealing that complex single-neuron responses often reflect the projection of simpler population-level dynamics onto individual cells \citep{vyas2020computation, saxena2019towards}. When viewed through population-level measurements, high-dimensional nonlinear neural systems often admit locally low-dimensional or approximately linear representations---an observation that opens the door to control-theoretic approaches for manipulating neural activity.

While observation of mesoscale dynamics has advanced rapidly, establishing causal links between these dynamics and behavioral outcomes remains challenging. Optogenetics has transformed the study of neural circuits by enabling temporally precise, cell-type-specific manipulation of neural activity \citep{deisseroth2011optogenetics, fenno2011development}. Most optogenetic experiments employ open-loop stimulation, in which predetermined patterns of light are delivered irrespective of the ongoing neural state. Although open-loop approaches have yielded fundamental insights into circuit function, they are inherently limited in their ability to account for the variability and state-dependence of neural responses \citep{grosenick2015closed}. Cortical activity fluctuates substantially from trial to trial due to ongoing spontaneous dynamics, arousal, and behavioral state \citep{harris2011cortical}, meaning that the same open-loop stimulus can produce widely varying neural and behavioral effects depending on the state of the brain at the time of stimulation.

Closed-loop paradigms, in which stimulation is continuously adjusted based on real-time measurements of neural activity, offer a principled framework for overcoming these limitations. From a control-theoretic perspective, feedback control enables a system to compensate for unmodeled disturbances and maintain a desired output despite variability in the underlying dynamics \citep{astrom1995pid}. In the context of neural manipulation, closing the loop provides two key advantages: it allows stimulation to adapt to the current brain state, improving the precision of neural modulation, and it enables the experimenter to drive neural activity toward specific target trajectories, creating opportunities for causal tests that are not possible with open-loop stimulation alone \citep{oshea2018importance}.

Several groups have demonstrated the feasibility and utility of closed-loop optogenetic control at the level of single neurons and local field potentials. \citet{newman2015optogenetic} introduced the ``optoclamp,'' a proportional--integral (PI) controller that locked the firing rate of individual neurons to stationary targets in cultured networks and in the anesthetized rat thalamus. \citet{bolus2018design} developed model-based design strategies for dynamic closed-loop optogenetic neurocontrol in vivo, and subsequently demonstrated state-space optimal feedback control of optogenetically driven single-neuron firing rates in the thalamus of awake, head-fixed mice using linear-quadratic optimal control \citep{bolus2021state}. At the network level, closed-loop optogenetic stimulation has been used to modulate oscillatory dynamics in brain slices and in non-human primates \citep{nicholson2018analogue, zaaimi2023closed}, including demonstrations of phase-dependent manipulation of seizure-like activity. \citet{kanta2019closed} used closed-loop optogenetic control of gamma oscillations in the amygdala to establish a causal link between endogenous gamma and memory consolidation. In a related line of work, \citet{clancy2014volitional} demonstrated that mice could learn to volitionally modulate optically recorded calcium signals through a neuroprosthetic interface, showing that closed-loop optical feedback can shape cortical activity patterns. A comprehensive review by \citet{grosenick2015closed} outlined the theoretical and practical foundations for closed-loop optogenetic control, emphasizing its potential for causal circuit dissection.

Despite these advances, most prior closed-loop optogenetic studies have operated at the level of single neurons or local field potentials, using electrophysiological recordings as the feedback signal. The application of closed-loop control to mesoscale cortical population dynamics---using wide-field calcium imaging as the readout---remains largely unexplored. This gap is significant because mesoscale dynamics reflect the integrated activity of large neuronal populations and are directly relevant to the kinds of distributed cortical computations that underlie perception and behavior. Moreover, recent technical developments have made it possible to perform simultaneous wide-field calcium imaging and spatially targeted optogenetic stimulation across the dorsal cortex of awake mice \citep{matveev2024simultaneous}, providing the experimental infrastructure needed for real-time, mesoscale closed-loop control.

Here, we demonstrate a closed-loop optogenetic control framework for regulating mesoscale cortical population activity. Using wide-field calcium imaging, we extract a real-time readout of regional cortical activity ($\Delta F/F$) and use this signal to drive a proportional--integral feedback controller that modulates targeted optogenetic stimulation of a neighboring cortical region. We first characterize the input--output relationship between optogenetic stimulation and trial-averaged cortical responses, showing that this relationship is well approximated by a linear mapping. We then implement a PI controller with feedforward compensation and demonstrate that closed-loop stimulation improves reference tracking accuracy and reduces trial-to-trial variability compared to open-loop stimulation across multiple experimental sessions and animals. Finally, we discuss the implications of these findings for the broader goal of using feedback-based perturbation paradigms to investigate and manipulate cortical population dynamics.

\bibliographystyle{apalike}
\bibliography{references}

\end{document}
